\clearpage
\VerdierPhysicalPageAnchor{0267}%
\begingroup
\small
\setlength{\parindent}{0pt}%
\setlength{\parskip}{0pt}%
\setlength{\emergencystretch}{2em}%
\newcommand{\VerdierBibRow}[3]{%
  \VerdierSemanticAnchor{#1}%
  \par\noindent\hangindent=3em\hangafter=1
  \makebox[3em][l]{[#2]}#3\par\smallskip}

\VerdierBibRow{bib-011}{11}{A. Grothendieck,
  \emph{Lefschetz 公式と L 関数の有理性}, Séminaire Bourbaki,
  exposé n° 279 (1964-1965).}

\VerdierBibRow{bib-012}{12}{J.-L. Verdier,
  \emph{エタール・コホモロジーにおける Lefschetz 不動点公式},
  Proceedings of a Conference on Local Fields,
  Nuffic Summer School held at Driebergen in 1966, pp. 199-214,
  T. A. Springer Ed., Springer-Verlag (1967).}

% Entry [13] is intentionally empty in the printed authority.
\VerdierBibRow{bib-013}{13}{}

\VerdierBibRow{bib-014}{14}{P. Berthelot、A. Grothendieck、L. Illusie,
  「交叉理論と Riemann-Roch 定理」 (S.G.A. 6),
  Séminaire de Géométrie Algébrique du Bois-Marie (1966-1967),
  Lecture Notes in Mathematics n° 225, Springer-Verlag (1971).}

\VerdierBibRow{bib-015}{15}{R. L. Knighten,「導来関手」,
  Thèse M.I.T. (1966).}

\VerdierBibRow{bib-016}{16}{P. Gabriel と M. Zisman,
  「分数計算とホモトピー論」, Ergebnisse der Mathematik, Band 35,
  Springer-Verlag (1967).}

\VerdierBibRow{bib-017}{17}{P. Freyd, \emph{アーベル圏における表現},
  Proceedings of the Conference on Categorical Algebra, La Jolla 1965,
  pp. 95-120, S. Eilenberg, D. K. Harrison, S. MacLane, H. Röhrl Eds.,
  Springer-Verlag (1966).}

\VerdierBibRow{bib-018}{18}{A. Grothendieck,
  \emph{ファイバー化された圏と降下},「エタール被覆と基本群」 (S.G.A. 1)
  Séminaire de Géométrie Algébrique du Bois-Marie (1960-1961), exposé n° VI,
  pp. 145-194, Lecture Notes in Mathematics n° 224, Springer-Verlag (1971).}

\VerdierBibRow{bib-019}{19}{J. Giraud, \emph{降下法}, Bull. Soc. Math.\ de
  France, Mémoire n° 2 (1964)}

\VerdierBibRow{bib-020}{20}{P. Gabriel, \emph{アーベル圏について},
  Bull. Soc. Math.\ de France, Tome 90, pp. 323-448 (1962).}

\VerdierBibRow{bib-021}{21}{P. Freyd, \emph{ホモトピー冪等射の分裂},
  Proceedings of the Conference on Categorical Algebra, La Jolla 1965,
  pp. 173-176, S. Eilenberg, D. K. Harrison, S. MacLane, H. Röhrl Eds.,
  Springer-Verlag (1966).}

\VerdierBibRow{bib-022}{22}{L. Gruson, \emph{アーベル完備化},
  Bull. Sciences Math., Tome 90, pp. 17-40 (1966).}

\VerdierBibRow{bib-023}{23}{R. Godement,「層の理論」,
  Actualités scientifiques et industrielles 1252, Hermann (1958).}

% Source-fidelity note: the printed authority reads "homologiqùe"; the
% localized title remains idiomatic.
\VerdierBibRow{bib-024}{24}{A. Grothendieck,
  \emph{ホモロジー代数のいくつかの点について}, Tôhoku Math. J., vol. 9,
  pp. 119-221 (1957).}
\endgroup
